\documentclass[../../main.tex]{subfiles}
\graphicspath{{\subfix{../../image/}}}

\begin{document}

\section{概率空间}\label{section:ch01-概率空间}

\subsection*{一、走向概率论公理化结构}

到本世纪初，概率论的各个领域已经得到了大量的成果，而且人们对概率论在其他基础学科和工程技术上的应用也产生了越来越大的兴趣。但是直到那时为止，关于概率论的一些基本概念——例如事件、概率等——却没有明确的定义。这是一个很大的矛盾，这个矛盾不仅可能导致贝特朗奇论那样的怪现象产生，而且也使许多人对概率的客观含义，甚至概率论结论的可应用性都产生了怀疑。因此可以说，到那时为止，概率论作为一个数学分支来说，还缺乏严格的理论基础，这就大大妨碍了它的进一步发展。

在概率论发展早期，所研究的随机现象比较简单，大部分可以归入古典概型。对这种模型，利用对称性概率的计算可以通过某种等可能性的假设进行，其结果也有相当明确的解释。这种成功使得人们试图通过给定某种等可能性来定义概率。于是，由 Laplace 提出的概率的古典定义在整个 $19$ 世纪被人们广泛接受。但是，这种定义的局限性很快也暴露了出来，它既要求试验的可能结果总数有限，又要求某种等可能性，所以它的适用范围极其有限。当把这个结果推广到有无限多种可能结果的场合，例如几何概率时，不但适用范围依然有限，而且出现了新问题。况且，用等可能性来定义概率还有循环定义之嫌。总之，对一般的随机现象明确地定义概率及其他基本概念，在那时成了一个突出的问题。

解决这个问题的时机也在不断成熟。首先是通过对概率论的两个最基本概念——事件与概率的长期研究，发觉事件的运算与集合的运算完全相似，概率与测度有相同性质，这个事实随着当时在实变函数论中关于 Lebesgue（$1875$—$1941$）测度和积分的研究以及一般抽象测度和积分理论的发展而日益明确起来。

另外，$19$ 世纪末以来，数学的各个分支广泛流行着一股公理化潮流，这个流派主张把最基本的假定公理化，其他结论则由它们经过演绎导出。

在这种背景下，$1933$ 年，前苏联数学家 Kolmogorov（$1903$—$1987$）提出了概率论公理化结构，这个结构综合了前人成果，明确定义了基本概念，使概率论成为严谨的数学分支，对近几十年来概率论的迅速发展起了积极作用。Kolmogorov 的这个理论已被普遍接受，本书的各部分就是在这个结构中展开的。

\subsection*{二、事件域}

在公理化结构中，事件不是最基本的概念，它通过更基本的概念——样本点来定义。样本点相应于随机试验的结果，我们已在前三节中进行过不少描述。在古典概型中，它们是可能结果全体，被用来定义一般的事件。在几何概率中，它们相应于区域 $\Omega$ 中的点，也被用来定义一般的事件。以后，我们把样本点 $\omega$ 看作抽象的点，它们的全体构成样本空间 $\Omega$。

正如 \S 2 中所做的那样，我们把事件 $A$ 定义为 $\Omega$ 的一个子集，它包含若干样本点，事件 $A$ 发生当且仅当 $A$ 所包含的样本点中有一个发生。

一般并不把 $\Omega$ 的一切子集都作为事件，因为这将对给定概率带来困难，譬如在几何概率中，若把不可测集也作为事件，将带来不可克服的困难。

另一方面，又必须把问题中感兴趣的事件都包括进来。例如若 $A$ 是事件，则应要求 $\overline{A}$ 也是事件；若 $A$ 与 $B$ 是事件，则 $A\cup B$ 及 $AB$ 也应是事件。当样本空间 $\Omega$ 由无限多个点构成时——在几何概率中就是如此——显然还得考虑可列个事件的并与交。此外，把 $\Omega$ 及 $\varnothing$ 作为事件有很大方便。

总之，我们若把事件的全体记为 $\mathscr{F}$，它是由 $\Omega$ 的一些子集构成的集类。而且为了使讨论便于进行，还得对 $\mathscr{F}$ 加上某些限制：

（i）$\Omega\in\mathscr{F}$；

（ii）若 $A\in\mathscr{F}$，则 $\overline{A}\in\mathscr{F}$；

（iii）若 $A_n\in\mathscr{F}$，$n=1,2,\cdots$，则 $\bigcup_{n=1}^{\infty}A_n\in\mathscr{F}$。

一般地，称空间 $\Omega$ 上满足上述三个要求的集类为 $\sigma$ \textbf{域}，亦称 $\sigma$ \textbf{代数}。

若 $\mathscr{F}$ 为 $\sigma$ 域，则由（i）及（ii）可得
\begin{align*}
\varnothing\in\mathscr{F}
\end{align*}
此外，若 $A_n\in\mathscr{F}$，$n=1,2,\cdots$，则
\begin{align*}
\bigcap_{n=1}^{\infty}A_n=\overline{\bigcup_{n=1}^{\infty}\overline{A_n}}\in\mathscr{F}
\end{align*}
\begin{align*}
\bigcup_{n=1}^{k}A_n=A_1\cup A_2\cup\cdots\cup A_k\cup\varnothing\cup\varnothing\cdots\in\mathscr{F}
\end{align*}
\begin{align*}
\bigcap_{n=1}^{k}A_n=A_1\cap A_2\cap\cdots\cap A_k\cap\Omega\cap\Omega\cdots\in\mathscr{F}
\end{align*}

据此定义，$\sigma$ 域对逆、并、交、差的可列次运算封闭，并且包含了 $\Omega$ 及 $\varnothing$。

\begin{definition}[事件域]\label{definition:ch01-事件域}
若 $\mathscr{F}$ 是由样本空间 $\Omega$ 的一些子集构成的一个 $\sigma$ 域，则称它为\textbf{事件域}（event field），$\mathscr{F}$ 中的元素称为\textbf{事件}，$\Omega$ 称为\textbf{必然事件}，$\varnothing$ 称为\textbf{不可能事件}。
\end{definition}

值得指出，按照这种定义，样本点并不一定是事件。

下面我们来举一些事件域的例子。

\begin{example}[平凡事件域]\label{example:ch01-平凡事件域}
$\mathscr{F}=\{\varnothing,\Omega\}$，不难验证 $\mathscr{F}$ 是一个 $\sigma$ 域，这时只有必然事件 $\Omega$ 与不可能事件 $\varnothing$ 是事件。
\end{example}

\begin{example}[二元事件域]\label{example:ch01-二元事件域}
$\mathscr{F}=\{\varnothing,A,\overline{A},\Omega\}$。这时 $\mathscr{F}$ 也是一个 $\sigma$ 域，$\Omega,A,\overline{A},\varnothing$ 是事件。
\end{example}

\begin{example}[有限集幂集]\label{example:ch01-有限集幂集}
$\Omega=\{\omega_1,\cdots,\omega_n\}$，$\mathscr{F}$ 由 $\Omega$ 的一切子集构成，这时它包含不可能事件 $\varnothing$，$\binom{n}{1}$ 个单点集，$\binom{n}{2}$ 个双点集，$\cdots$，$\binom{n}{n-1}$ 个 $n-1$ 点集，还有必然事件 $\Omega$，因此计有
\begin{align*}
\binom{n}{0}+\binom{n}{1}+\binom{n}{2}+\cdots+\binom{n}{n-1}+\binom{n}{n}=2^n
\end{align*}
个元素。不难验证，$\mathscr{F}$ 是一个 $\sigma$ 域。
\end{example}

\begin{example}[一般幂集]\label{example:ch01-一般幂集}
对于一般的 $\Omega$，若 $\mathscr{F}$ 由 $\Omega$ 的一切子集构成，可以验证 $\mathscr{F}$ 是一个 $\sigma$ 域。
\end{example}

从上面几个例子中看到，事件域可以很简单，也可以选得十分复杂，这就需要我们根据问题的不同要求来选择适当的事件域。

表面看，当 $\Omega$ 确定后，把事件域 $\mathscr{F}$ 选得越大，能处理的事件越多，就越方便。但是概率论最关心的毕竟是概率，过大的事件域对概率的给定带来困难，并不可取。不过，如果定义概率没有困难，那么，事件域当然可以尽量选大。因此对有限样本空间和离散样本空间，以后我们将看到，通常都取 $\Omega$ 的一切子集作为事件域。

对一个试验，当 $\Omega$ 给定后，总有些子集必须作为事件处理，但它们未必能满足 $\sigma$ 域的要求，该怎么办？

下面证明：若给定 $\Omega$ 的一个非空集类 $\mathscr{G}$，必存在唯一的一个 $\Omega$ 上的 $\sigma$ 域 $\mathfrak{m}(\mathscr{G})$，具有如下两个性质：（1）包含 $\mathscr{G}$；（2）若有其他 $\sigma$ 域包含 $\mathscr{G}$，则必包含 $\mathfrak{m}(\mathscr{G})$。这个 $\mathfrak{m}(\mathscr{G})$ 称为包含 $\mathscr{G}$ 的\textbf{最小 $\sigma$ 域}，亦称由 $\mathscr{G}$ \textbf{产生的 $\sigma$ 域}。

先证明必存在包含 $\mathscr{G}$ 的 $\sigma$ 域，显然由 $\Omega$ 的一切子集构成的集类包含了 $\mathscr{G}$，由例~\ref{example:ch01-一般幂集} 知此集类是一个 $\sigma$ 域，因此至少存在 $\Omega$ 上的一个 $\sigma$ 域 $\mathfrak{m}$，有 $\mathfrak{m}\supset\mathscr{G}$。

现在只要取 $\Omega$ 上一切包含 $\mathscr{G}$ 的 $\sigma$ 域之交作为 $\mathfrak{m}(\mathscr{G})$，则它是具有上述两个性质的 $\sigma$ 域。这点作为习题留给读者自行验证。

因此，从必须作为事件处理的子集出发，通过添加其他子集，必能得到 $\Omega$ 上的 $\sigma$ 域。不过，上述最小 $\sigma$ 域因便于给定概率而受重视。

按照这种观点，例~\ref{example:ch01-平凡事件域} 是只把不可能事件 $\varnothing$ 及必然事件 $\Omega$ 看作事件的平凡事件域，而例~\ref{example:ch01-二元事件域} 是由事件 $A$ 产生的事件域。

当 $\Omega$ 有限或可列，如果要求每一个样本点都是事件，则包含它的最小 $\sigma$ 域就是 $\Omega$ 的一切子集（如例~\ref{example:ch01-有限集幂集}）。因此在这两种场合，事件域的选取实际上没有困难。

真正要关心的是样本空间为一维或 $n$ 维 Euclid 空间的场合。这里的许多结果是由 Borel（$1871$—$1956$）建立的。

至于如何在函数空间上定义事件域则超出本书的讨论范围。

下面介绍两个非常有用的 $\sigma$ 域。

\textbf{[一维 Borel 点集]}~~以后我们将以 $\mathbf{R}^1$ 记数直线或实数全体，并称由一切形为 $[a,b)$ 的有界左闭右开区间构成的集类所产生的 $\sigma$ 域为\textbf{一维 Borel $\sigma$ 域}，记之为 $\mathscr{B}_1$，称 $\mathscr{B}_1$ 中的集为\textbf{一维 Borel 点集}。

若 $x,y$ 为任意实数，由于
\begin{align*}
\{x\}=\bigcap_{n=1}^{\infty}\left[x,x+\frac{1}{n}\right)
\end{align*}
\begin{align*}
(x,y)=[x,y)-\{x\}
\end{align*}
\begin{align*}
[x,y]=[x,y)+\{y\}
\end{align*}
\begin{align*}
(x,y]=[x,y)+\{y\}-\{x\}
\end{align*}
因此 $\mathscr{B}_1$ 中包含一切开区间，闭区间，单个实数，可列个实数，以及由它们经可列次逆、并、交运算而得出的集合。这是相当大的一个集类，足够把实际问题中感兴趣的点集都包括在内。

显然，若不从左闭右开区间 $[a,b)$ 出发，而从 $(a,b]$ 或 $(a,b)$，或 $[a,b]$，甚至 $(-\infty,x)$ 出发，都将产生同一个 $\sigma$ 域。

\textbf{[$n$ 维 Borel 点集]}~~以 $\mathbf{R}^n$ 记 $n$ 维 Euclid 空间，可以类似地定义 $n$ 维 Borel 点集，它们是由一切 $n$ 维矩形产生的 $n$ 维 Borel $\sigma$ 域 $\mathscr{B}_n$ 中的集合，也可以把 $\mathbf{R}^n$ 中我们感兴趣的点集都包括在内。

\subsection*{三、概率}

在公理化结构中，概率是针对事件定义的，即对应于事件域 $\mathscr{F}$ 中的每一个元素 $A$ 有一个实数 $P(A)$ 与之对应，一般把这种从集合到实数的映射称为集合函数。因此，概率是定义在事件域 $\mathscr{F}$ 上的一个集合函数。此外，在公理化结构中只规定概率应满足的性质，而不具体给出它的计算公式或计算方法。

概率应有什么性质呢？

因为概率通过频率稳定性与随机试验相联系，因此我们自然想到概率应有与频率类似的性质。关于频率的性质，我们已在 \S 1 中总结为非负性，规范性以及有限可加性。

在古典概型中，概率是通过有利场合数与可能结果总数之比来定义的，它同样具有这三个性质。

在几何概率中，情况也类似，但有一点不同，就是它要求对可列个不相容事件之和有可加性，即可列可加性。

在一般场合，处理可列个事件之和是完全必要的，因此保留这种可列可加性要求看来是合理的。

综上所述，我们如下定义概率。

\begin{definition}[概率]\label{definition:ch01-概率}
定义在事件域 $\mathscr{F}$ 上的一个集合函数 $P$ 称为\textbf{概率}，如果它满足如下三个要求：

（i）$P(A)\geqslant 0$，对一切 $A\in\mathscr{F}$；

（ii）$P(\Omega)=1$；

（iii）若 $A_i\in\mathscr{F}$，$i=1,2,\cdots$ 且两两互不相容，则
\begin{align}
P\left(\sum_{i=1}^{\infty}A_i\right)=\sum_{i=1}^{\infty}P(A_i)\label{align:ch01-概率可列可加}
\end{align}
性质（i）称为\textbf{非负性}，性质（ii）称为\textbf{规范性}，性质（iii）称为\textbf{可列可加性}或\textbf{完全可加性}。
\end{definition}

利用概率的基本性质（i），（ii），（iii）可以推出概率的另外一些重要性质。

\begin{theorem}[不可能事件概率]\label{theorem:ch01-不可能事件概率}
不可能事件的概率为 $0$，即 $P(\varnothing)=0$。
\end{theorem}

\begin{proof}
因为
\begin{align*}
\Omega=\Omega+\varnothing+\cdots
\end{align*}
所以
\begin{align*}
P(\Omega)=P(\Omega)+P(\varnothing)+\cdots
\end{align*}
因此 $P(\varnothing)=0$。
\end{proof}

\begin{theorem}[有限可加性]\label{theorem:ch01-有限可加性}
概率具有有限可加性。即若 $A_iA_j=\varnothing$（$i\neq j$），则
\begin{align}
P(A_1+A_2+\cdots+A_n)=P(A_1)+P(A_2)+\cdots+P(A_n)\label{align:ch01-有限可加性}
\end{align}
\end{theorem}

\begin{proof}
因为
\begin{align*}
A_1+A_2+\cdots+A_n=A_1+A_2+\cdots+A_n+\varnothing+\varnothing+\cdots
\end{align*}
由可列可加性及定理~\ref{theorem:ch01-不可能事件概率} 得
\begin{align*}
P(A_1+A_2+\cdots+A_n)=P(A_1)+P(A_2)+\cdots+P(A_n)
\end{align*}
\end{proof}

\begin{theorem}[逆事件概率]\label{theorem:ch01-逆事件概率一般}
对任何事件 $A$ 有
\begin{align}
P(\overline{A})=1-P(A)\label{align:ch01-逆事件概率一般}
\end{align}
\end{theorem}

\begin{proof}
因 $A\cup\overline{A}=\Omega$，$A\overline{A}=\varnothing$，故
\begin{align*}
1=P(\Omega)=P(A\cup\overline{A})=P(A)+P(\overline{A})
\end{align*}
\end{proof}

\begin{theorem}[差事件概率]\label{theorem:ch01-差事件概率}
如果 $A\supset B$，则
\begin{align}
P(A-B)=P(A)-P(B)\label{align:ch01-差事件概率}
\end{align}
\end{theorem}

\begin{proof}
因为 $A=AB+A\overline{B}$，故
\begin{align*}
P(A-B)=P(A\overline{B})=P(A)-P(AB)
\end{align*}
特别当 $A\supset B$ 时，得到式~\eqref{align:ch01-差事件概率}。
\end{proof}

\begin{corollary}[单调性]\label{corollary:ch01-单调性}
如果 $A\supset B$，则 $P(A)\geqslant P(B)$。
\end{corollary}

由此即知，对任意事件 $A$，有 $P(A)\leqslant 1$，再注意到概率的非负性，所以成立
\begin{align}
0\leqslant P(A)\leqslant 1\label{align:ch01-概率界}
\end{align}

\begin{theorem}[加法公式]\label{theorem:ch01-加法公式}
\begin{align}
P(A\cup B)=P(A)+P(B)-P(AB)\label{align:ch01-加法公式}
\end{align}
\end{theorem}

\begin{proof}
因 $A\cup B=A\cup(B-AB)$，而且 $A\cap(B-AB)=\varnothing$，故 $P(A\cup B)=P(A)+P(B-AB)$，又 $AB\subset B$，于是由定理~\ref{theorem:ch01-差事件概率} 得到
\begin{align*}
P(A\cup B)=P(A)+P(B)-P(AB)
\end{align*}
\end{proof}

公式~\eqref{align:ch01-加法公式} 称为概率的\textbf{加法公式}，在古典概型场合，我们曾直接加以证明，如今则作为可加性的直接推论，说明在一般场合也成立。由它可得如下推论：

\begin{corollary}[Boole 不等式]\label{corollary:ch01-boole}
\begin{align}
P(A\cup B)\leqslant P(A)+P(B)\label{align:ch01-boole}
\end{align}
\end{corollary}

\begin{corollary}[Bonferroni 不等式]\label{corollary:ch01-bonferroni}
\begin{align}
P(AB)\geqslant P(A)+P(B)-1\label{align:ch01-bonferroni}
\end{align}
\end{corollary}

利用归纳法不难把这两个不等式推广到 $n$ 个事件的场合。
\begin{align*}
P(A_1\cup A_2\cup\cdots\cup A_n)\leqslant P(A_1)+P(A_2)+\cdots+P(A_n)
\end{align*}
\begin{align*}
P(A_1A_2\cdots A_n)\geqslant P(A_1)+P(A_2)+\cdots+P(A_n)-(n-1)
\end{align*}

\begin{theorem}[一般加法公式]\label{theorem:ch01-一般加法公式}
若 $A_1,A_2,\cdots,A_n$ 为 $n$ 个事件，则
\begin{align}
P(A_1\cup A_2\cup\cdots\cup A_n)&=\sum_{i=1,\cdots,n}P(A_i)-\sum_{\substack{i<j\\i,j=1,\cdots,n}}P(A_iA_j)\notag\\
&\quad+\sum_{\substack{i<j<k\\i,j,k=1,\cdots,n}}P(A_iA_jA_k)-\cdots+(-1)^{n-1}P(A_1A_2\cdots A_n)\label{align:ch01-一般加法公式}
\end{align}
\end{theorem}

公式~\eqref{align:ch01-一般加法公式} 可以用数学归纳法证明，留给读者作为习题。

当然，概率还有其他性质，不过上述列出的各条是最重要的。它们是以后讨论中经常要用到的基本结果，务必牢记。

值得提醒的是概率的这些重要性质的推导实质上只用到非负性、规范性和有限可加性。

\textbf{[概率计算的公式法]}~~利用上面推导的公式来作概率计算，常能使解题思路清晰，计算便捷。

\begin{example}[最大车牌号]\label{example:ch01-最大车牌号}
某城有 $N$ 辆卡车，车牌号从 $1$ 到 $N$，有一个外地人到该城去，把遇到的 $n$ 辆车子的牌号抄下（可能重复抄到某些车牌号），求抄到的最大号码正好为 $k$ 的概率（$1\leqslant k\leqslant N$）。

\textbf{解}~~可以看作古典概型问题，即设每辆卡车被遇到的机会相同。若以 $A_k$ 记抄到的最大号码为 $k$ 这一事件，又以 $B_k$ 记抄到的最大号码不超过 $k$ 这一事件，则明显有 $A_k=B_k-B_{k-1}$，而且 $B_k\supset B_{k-1}$，所以由定理~\ref{theorem:ch01-差事件概率} 知 $P(A_k)=P(B_k)-P(B_{k-1})$，而由直接计算可得 $P(B_k)=\frac{k^n}{N^n}$，因此最后得到
\begin{align*}
P(A_k)=\frac{k^n-(k-1)^n}{N^n}
\end{align*}

这种方法曾在第二次世界大战中被盟军用来估计敌方的军火生产能力，从被击毁的战车上的出厂号码推测其生产批量，得到相当精确的有用情报。
\end{example}

\begin{example}[匹配问题]\label{example:ch01-匹配问题}
某人写好 $n$ 封信，又写好 $n$ 只信封，然后在黑暗中把每封信放入一只信封中，试求至少有一封信放对的概率。

\textbf{解}~~若以 $A_i$ 记第 $i$ 封信与信封符合，则所求事件为 $A_1\cup A_2\cup\cdots\cup A_n$，所以可以用一般加法公式，不难求得
\begin{align*}
P(A_i)=\frac{(n-1)!}{n!}=\frac{1}{n}
\end{align*}
\begin{align*}
P(A_iA_j)=\frac{(n-2)!}{n!}=\frac{1}{n(n-1)}
\end{align*}
\begin{align*}
P(A_iA_jA_k)=\frac{1}{n(n-1)(n-2)},\cdots,P(A_1A_2\cdots A_n)=\frac{1}{n!}
\end{align*}
因此
\begin{align*}
P(A_1\cup A_2\cup\cdots\cup A_n)&=\binom{n}{1}\frac{1}{n}-\binom{n}{2}\frac{1}{n(n-1)}\\
&\quad+\binom{n}{3}\frac{1}{n(n-1)(n-2)}-\cdots+(-1)^{n-1}\frac{1}{n!}\\
&=1-\frac{1}{2!}+\frac{1}{3!}-\cdots+(-1)^{n-1}\frac{1}{n!}
\end{align*}
这个问题若直接计算有利场合数，无疑是十分复杂的。
\end{example}

\subsection*{四、可列可加性与连续性}

下面我们对可列可加性作进一步讨论。从定理~\ref{theorem:ch01-有限可加性} 知道，由可列可加性可以推出有限可加性，但是一般来讲，由有限可加性并不能推出可列可加性。

事实上，若 $A_i\in\mathscr{F}$，$i=1,2,\cdots$ 且两两互不相容，则由概率的有限可加性只能推出~\eqref{align:ch01-有限可加性} 式成立，即
\begin{align}
P\left(\sum_{i=1}^{n}A_i\right)=\sum_{i=1}^{n}P(A_i)\label{align:ch01-有限和}
\end{align}
这个等式的左边对任意 $n$ 都不超过 $1$，因此右边的正项级数收敛。这样应有
\begin{align}
\lim_{n\to\infty}P\left(\sum_{i=1}^{n}A_i\right)=\lim_{n\to\infty}\sum_{i=1}^{n}P(A_i)=\sum_{i=1}^{\infty}P(A_i)\label{align:ch01-级数收敛}
\end{align}
与~\eqref{align:ch01-概率可列可加} 式比较一下就可以知道，为了具有可列可加性，还需要下式成立：
\begin{align}
\lim_{n\to\infty}P\left(\sum_{i=1}^{n}A_i\right)=P\left(\sum_{i=1}^{\infty}A_i\right)\label{align:ch01-交换极限}
\end{align}
或者写成更富有启发性的等式：
\begin{align}
\lim_{n\to\infty}P\left(\sum_{i=1}^{n}A_i\right)=P\left(\lim_{n\to\infty}\sum_{i=1}^{n}A_i\right)\label{align:ch01-交换极限启发性}
\end{align}
即要允许把极限号移到概率号里面去，这就提出了一个新要求。

现在就来考察这个新要求。若记
\begin{align*}
S_n=\sum_{i=1}^{n}A_i
\end{align*}
则 $S_n\in\mathscr{F}$，$n=1,2,\cdots$ 而且 $S_n\subset S_{n+1}$，即 $S_n$ 是 $\mathscr{F}$ 中一个单调不减的集序列，这时可改写~\eqref{align:ch01-交换极限启发性} 式成为下式：
\begin{align}
\lim_{n\to\infty}P(S_n)=P\left(\lim_{n\to\infty}S_n\right)\label{align:ch01-下连续}
\end{align}

一般，对于 $\mathscr{F}$ 上的集合函数，若它对 $\mathscr{F}$ 中任何一个单调不减的集序列 $\{S_n\}$ 均成立~\eqref{align:ch01-下连续} 式，则我们称它是\textbf{下连续的}。因此前面的推导表明，为保证概率的可列可加性成立，除要求它具有有限可加性外，还要求它是下连续的。

下面的定理明确地阐明了这三个概念之间的关系。

\begin{theorem}[可列可加性充要条件]\label{theorem:ch01-可列可加充要条件}
若 $P$ 是 $\mathscr{F}$ 上满足 $P(\Omega)=1$ 的非负集合函数，则它具有可列可加性的充要条件为

（i）它是有限可加的；

（ii）它是下连续的。
\end{theorem}

\begin{proof}
\textbf{充分性}~~只要重新考察前面的推导过程。实际上，若沿用有关记号，则有限可加性保证了~\eqref{align:ch01-有限和} 式的成立，而下连续性保证了~\eqref{align:ch01-交换极限启发性} 式即~\eqref{align:ch01-交换极限} 式的成立，通过~\eqref{align:ch01-级数收敛} 式立刻得到可列可加性。

\textbf{必要性}~~其中（i）我们早已建立。为证（ii），设 $\{S_n\}$ 是 $\mathscr{F}$ 中一个单调不减的集序列，那么
\begin{align*}
\bigcup_{i=1}^{\infty}S_i=\lim_{n\to\infty}S_n
\end{align*}
若定义 $S_0=\varnothing$，则
\begin{align*}
\bigcup_{i=1}^{\infty}S_i=\sum_{i=1}^{\infty}(S_i-S_{i-1})
\end{align*}
这里的 $(S_i-S_{i-1})$，$i=1,2,\cdots$，由于 $S_i$ 的单调性，显然两两互不相容，因此由可列可加性得
\begin{align*}
P\left(\bigcup_{i=1}^{\infty}S_i\right)=\sum_{i=1}^{\infty}P(S_i-S_{i-1})=\lim_{n\to\infty}\sum_{i=1}^{n}P(S_i-S_{i-1})
\end{align*}
但是
\begin{align*}
\sum_{i=1}^{n}P(S_i-S_{i-1})=P\left(\sum_{i=1}^{n}(S_i-S_{i-1})\right)=P(S_n)
\end{align*}
因此
\begin{align*}
P\left(\lim_{n\to\infty}S_n\right)=\lim_{n\to\infty}P(S_n)
\end{align*}
这就证得了 $P$ 的下连续性。
\end{proof}

\begin{corollary}\label{corollary:ch01-下连续}
概率是下连续的。
\end{corollary}

\begin{corollary}[上连续]\label{corollary:ch01-上连续}
概率是上连续的，即若 $B_i\in\mathscr{F}$，而且 $B_i\supset B_{i+1}$，$i=1,2,\cdots$，则
\begin{align}
\lim_{n\to\infty}P(B_n)=P\left(\lim_{n\to\infty}B_n\right)\label{align:ch01-上连续}
\end{align}
\end{corollary}

\begin{proof}
记 $S_i=\overline{B_i}$，则 $\{S_i\}$ 是单调不减的，由推论~\ref{corollary:ch01-下连续} 可得
\begin{align*}
\lim_{n\to\infty}P(S_n)=P\left(\lim_{n\to\infty}S_n\right)=P\left(\bigcup_{i=1}^{\infty}S_i\right)=P\left(\overline{\bigcap_{i=1}^{\infty}B_i}\right)
\end{align*}
因此由~\eqref{align:ch01-逆事件概率一般} 式，可知
\begin{align*}
1-\lim_{n\to\infty}P(B_n)=1-P\left(\bigcap_{i=1}^{\infty}B_i\right)
\end{align*}
注意到 $\lim_{n\to\infty}B_n=\bigcap_{i=1}^{\infty}B_i$，即得式~\eqref{align:ch01-上连续}。
\end{proof}

\begin{corollary}\label{corollary:ch01-次可加}
\begin{align*}
P\left(\bigcup_{i=1}^{\infty}A_i\right)\leqslant\sum_{i=1}^{\infty}P(A_i)
\end{align*}
\end{corollary}

\subsection*{五、概率空间}

在 Kolmogorov 的概率论公理化结构中，称三元总体 $(\Omega,\mathscr{F},P)$ 为\textbf{概率空间}，其中 $\Omega$ 是样本空间，$\mathscr{F}$ 是事件域，$P$ 是概率，它们都认为是预先给定的，并以此作为出发点讨论种种问题。至于实际问题中，如何选定 $\Omega$，怎样构造 $\mathscr{F}$，怎样给定 $P$，则要视具体情况而定。

下面讨论几个具体例子。

\begin{example}[有限概率空间]\label{example:ch01-有限概率空间}
设 $\Omega$ 中只有 $n$ 个点，这类概率空间在 \S 2 中已有讨论，古典概型即为其特例。在这种场合，一般可取 $\mathscr{F}$ 为 $\Omega$ 的所有子集全体，这仍是一个有限的集合，元素总数为 $2^n$，它满足事件域的三个要求，而且样本点（看作一个单点集）是事件。至于概率，只要对样本点 $\omega_i$，$i=1,2,\cdots,n$，给定满足
\begin{align*}
p(\omega_i)\geqslant 0,\qquad i=1,2,\cdots,n
\end{align*}
\begin{align*}
p(\omega_1)+p(\omega_2)+\cdots+p(\omega_n)=1
\end{align*}
的一组数 $p(\omega_1),p(\omega_2),\cdots,p(\omega_n)$，那么，若 $A$ 是 $\mathscr{F}$ 中元素，包含样本点 $\omega_{i_1},\cdots,\omega_{i_k}$，则由概率的可加性，自然应令
\begin{align*}
P(A)=p(\omega_{i_1})+\cdots+p(\omega_{i_k})
\end{align*}
这就给定了事件 $A$ 的概率，从而构成了概率空间 $(\Omega,\mathscr{F},P)$。

这时显然有 $P(\{\omega_i\})=p(\omega_i)$，$i=1,2,\cdots,n$。

从这个例子中看到下面两点。

（1）选定了 $(\Omega,\mathscr{F})$ 之后，对于事件概率的给定还有相当大的灵活性，这表现在 $p(\omega_i)$ 的选取上。因为只有这样，才能用概率空间来描述不同的随机现象。例如在投一次硬币的试验中，$\Omega$ 总是由出正面（$\omega_1$）及出反面（$\omega_2$）两个样本点构成。对于均匀的硬币，可以假定它出正面及反面的概率均为 $\frac{1}{2}$；但对于很不均匀的硬币，例如出正面可能性大得多的硬币，则必须给定另外的概率，而这只须适当给定 $p(\omega_1)$ 就可以了。

（2）一旦诸 $p(\omega_i)$ 给定后，事件 $A$ 的概率并不能任意给定，即在事件域中，各事件的概率有一定关系，给定概率时必须满足这些关系。
\end{example}

\begin{example}[离散概率空间]\label{example:ch01-离散概率空间}
$\Omega$ 由可列个点构成：$\Omega=\{\omega_1,\omega_2,\cdots\}$，这时 $\mathscr{F}$ 还是可以选为 $\Omega$ 的子集全体，它满足事件域的三个要求。这时样本点也是事件。为给定概率，可选择可列个非负的数 $p_i$，$i=1,2,\cdots$ 满足
\begin{align*}
\sum_{i=1}^{\infty}p_i=1
\end{align*}
分别作为样本点 $\omega_i$ 的概率，而一般事件 $A\in\mathscr{F}$ 的概率，则必须取为它所含的样本点的概率之和。
\end{example}

\begin{example}[实数上的概率空间]\label{example:ch01-实数概率空间}
若 $\Omega=\mathbf{R}^1$，即样本空间由全体实数构成，这时 $\mathscr{F}$ 不能取为 $\Omega$ 的一切子集，因为这个集类太大，无法在其上定义概率。这时通常取 $\mathscr{F}$ 为直线上 Borel 点集全体 $\mathscr{B}_1$，这是相当大的一个集类，可以把实际问题中所有感兴趣的点集都包括在内。另一方面在 Borel $\sigma$ 域上定义概率相当方便，这只要对左闭右开区间给定概率即可。这些我们将在第三章中作深入讨论。
\end{example}

顺便指出，若 $\Omega$ 不是 $\mathbf{R}^1$，而是它的一部分，也可类似处理。譬如 $\Omega$ 为一个区间，这时 $\mathscr{F}$ 可取为该区间上的 Borel 点集全体：它们通过直线上 Borel 点集与该区间之交而得到。

\begin{example}[高维概率空间]\label{example:ch01-高维概率空间}
若 $\Omega=\mathbf{R}^n$（或 $\mathbf{R}^n$ 的一部分），这时可类似于一维场合取 $n$ 维 Euclid 空间中的 Borel 点集全体 $\mathscr{B}_n$ 作为事件域 $\mathscr{F}$，在第三章 \S 2 中将对这种场合进行深入讨论。
\end{example}

\end{document}